\subsubsection{Subproceso 2: Selección y Registro de Cursos}


\textbf{Personas:}
\begin{itemize}
    \item \textbf{Estudiante Regular:} Usuario activo del proceso, responsable de tomar decisiones de selección e interactuar directamente con la plataforma.
    \item \textbf{Personal de la Comisión de Matrícula y Departamentos Académicos:} Encargados de la parametrización previa de la oferta académica y del soporte técnico-administrativo durante el periodo.
\end{itemize}

\textbf{Tecnología:}
\begin{itemize}
    \item \textbf{Intranet DIRCE (\texttt{https://dirce.uni.edu.pe/}):} Plataforma web institucional a través de la cual el estudiante visualiza la oferta, selecciona asignaturas y registra su propuesta de matrícula.
    \item \textbf{Sistema Integrado de Gestión Académica (SIGA-ORCE):} Sistema de información centralizado que procesa las transacciones concurrentes, controla los cupos disponibles en tiempo real y almacena la Pre-Ficha registrada.
    \item \textbf{Dispositivo de acceso con conexión a internet:} Computadora personal, portátil o dispositivo móvil con navegador web compatible utilizado por el estudiante para ingresar al portal.
\end{itemize}

\textbf{Documentos y referencias:}
\begin{itemize}
    \item \textbf{Malla Curricular Digital:} Mapeo automatizado del Plan de Estudios con la matriz de precedencias (prerrequisitos) que condiciona la visibilidad de asignaturas en la interfaz.
    \item \textbf{Tabla de límites de créditos (RR-0570-2022):} Parámetro normativo configurado en el sistema que determina el tope crediticio individual según el promedio ponderado del estudiante.
    \item \textbf{Cronograma de turnos de matrícula:} Documento publicado por la Facultad que indica la fecha y franja horaria habilitada para cada estudiante.
\end{itemize}



